\chapter{Requirements}

This chapter defines the main functional requirements, non-functional requirements, software design concepts, and design principles for the cardiopulmonary sound separation system. The system is treated as a prototype and proof of concept for a software design assignment and Final Year Project context. It is designed to demonstrate a realistic local web application that can upload mixed cardiopulmonary audio, run NeoSSNet-based separation, store metadata, and present results clearly to students, lecturers, researchers, or demonstration users. The requirements do not claim clinical validation; they focus on software behaviour, maintainability, and demonstration readiness.

\section{Functional Requirements}

Functional requirements describe the main behaviours that the system must provide to the user and to the backend workflow.

\begin{table}[H]
  \centering
  \small
  \caption{Functional Requirements}
  \label{tab:functional-requirements}
  \begin{tabularx}{\textwidth}{p{0.12\textwidth} p{0.25\textwidth} Y}
    \toprule
    \textbf{ID} & \textbf{Requirement} & \textbf{Description} \\
    \midrule
    FR-01 & Upload mixed WAV audio & The user shall be able to select and upload a mixed cardiopulmonary WAV file through the web interface. \\
    FR-02 & Validate audio file & The backend shall check that the uploaded file is a valid supported audio file before saving it for processing. \\
    FR-03 & Run cardiopulmonary sound separation & The backend shall run the separation workflow using NeoSSNet to separate the mixed input into heart and lung sound outputs. \\
    FR-04 & Generate heart and lung output files & The system shall save separated heart and lung audio files in the configured output folders. \\
    FR-05 & Preview separated audio & The web interface shall allow the user to preview the original audio and the separated heart and lung outputs where browser playback is supported. \\
    FR-06 & Download separated audio & The user shall be able to download the generated heart and lung output files for later review or demonstration. \\
    FR-07 & View processing history & The system shall show previous uploads, jobs, statuses, and result links so that users can review earlier processing activity. \\
    FR-08 & Store logs and metadata & The system shall store upload metadata, job records, result paths, metrics where available, and important processing logs in SQLite. \\
    \bottomrule
  \end{tabularx}
\end{table}

These functional requirements support the main demonstration flow of the project. A user uploads a mixed sound, the backend validates and processes it, the system stores the generated files, and the browser presents previews, downloads, and history. This keeps the project focused on a working end-to-end software system rather than only a standalone machine learning script.

\section{Non-Functional Requirements}

Non-functional requirements describe the quality attributes expected from the prototype. These qualities are important because the system should be understandable, maintainable, and reliable enough for local testing and presentation.

\begin{table}[H]
  \centering
  \small
  \caption{Non-Functional Requirements}
  \label{tab:non-functional-requirements}
  \begin{tabularx}{\textwidth}{p{0.24\textwidth} Y}
    \toprule
    \textbf{Quality Attribute} & \textbf{Requirement} \\
    \midrule
    Usability & The interface should make the upload, separation, preview, download, and history flow clear for a student or examiner during a demonstration. \\
    Maintainability & The code should be organised into clear frontend, router, service, machine learning, database, and storage responsibilities so that future changes are easier to make. \\
    Reliability & The backend should handle invalid files, failed separation jobs, missing outputs, and database errors with clear status messages and logs. \\
    Modularity & Each major part of the system should be replaceable or testable with minimal impact on unrelated components. \\
    Performance & The system should process short demonstration WAV files within a reasonable time on a local machine, while avoiding unnecessary database storage of large audio blobs. \\
    Data integrity & SQLite records should consistently link uploaded audio, model records, separation jobs, results, metrics, and logs. Failed jobs should keep useful error information. \\
    Portability and local execution & The system should run locally using Python, FastAPI, SQLite, PyTorch, and local folders without requiring cloud services or complex deployment infrastructure. \\
    \bottomrule
  \end{tabularx}
\end{table}

The non-functional requirements guide implementation decisions. For example, local execution supports a practical student project setup, while data integrity and logging support traceability when explaining how a mixed input becomes separated heart and lung outputs.

\section{Software Design Concepts}

\subsection{Abstraction}

Abstraction is used to hide detailed processing steps behind simpler operations. In this project, the separation route does not need to know how model records are selected, how NeoSSNet is called, or how heart and lung output paths are written. It can call a service-level operation such as \texttt{separate\_uploaded\_audio} and treat the workflow as one meaningful software action.

The machine learning boundary also uses abstraction. \texttt{SeparationAlgorithm} represents the expected behaviour of a separation model, while \texttt{NeoSSNetStrategy} provides the current real implementation. This allows the service layer to express the workflow in terms of cardiopulmonary sound separation rather than in terms of one concrete model class.

\subsection{Modularity}

Modularity divides the system into smaller parts with clear responsibilities. The frontend handles upload controls, audio preview, download links, and history display. FastAPI routers handle HTTP requests. Services coordinate workflows. The ML layer performs real separation. SQLite models store structured records, and local folders store uploaded and generated WAV files.

This modular structure is useful for the prototype because a change in one area does not require rewriting the entire system. For example, improving the frontend model dropdown should not require changing NeoSSNet inference, and changing output folder naming should not require changing the database schema.

\subsection{Separation of Concerns}

Separation of concerns prevents unrelated responsibilities from being mixed together. The upload route receives files, the storage service manages local file placement, the model service reads available model records, the separation service coordinates processing, and the ML strategy runs inference. This keeps web request handling, file management, model selection, and audio separation as separate concerns.

The same idea applies to data storage. Runtime upload folders, generated output folders, model checkpoints, dataset folders, and SQLite metadata serve different purposes and remain separate. This design avoids training from runtime upload folders and avoids storing large WAV content directly in SQLite.

\subsection{Information Hiding}

Information hiding means that internal implementation details are kept inside the module or class responsible for them. File validation and path rules should stay inside upload or storage-related code. Model selection rules should stay inside \texttt{ModelService}. Algorithm construction should stay inside \texttt{SeparationAlgorithmFactory}. NeoSSNet inference details should stay inside the ML strategy and inference function.

This protects the rest of the system from unnecessary implementation knowledge. For example, the route does not need to know where the checkpoint file is stored or how the strategy calls the underlying inference function; it only needs the job result returned by the service.

\subsection{Low Coupling and High Cohesion}

Low coupling means that components should depend on each other only where necessary. In this project, \texttt{SeparationService} depends on the \texttt{SeparationAlgorithm} abstraction rather than depending directly on NeoSSNet-specific workflow logic. This reduces the impact of future model changes because the upload, result, history, and download routes can remain stable.

High cohesion means that each component should focus on related tasks. \texttt{StorageService} is cohesive when it handles file storage concerns, \texttt{ModelService} is cohesive when it handles model lookup, and \texttt{NeoSSNetStrategy} is cohesive when it only adapts NeoSSNet inference to the shared algorithm interface. Cohesion helps the system remain easy to explain and test.

\section{Design Principles}

\subsection{Single Responsibility Principle}

The Single Responsibility Principle supports the assignment of one main reason for change to each component. For example, \texttt{StorageService} should change when file storage rules change, not when separation model logic changes. This keeps file handling separate from ML inference and database workflow decisions.

\begin{lstlisting}[caption={Single Responsibility Principle Example}]
# Single Responsibility Principle:
# StorageService handles storage concerns only.
class StorageService:
    def save_uploaded_wav(self, upload_file) -> str:
        stored_path = build_upload_path(upload_file.filename)
        write_file(upload_file, stored_path)
        return stored_path


# Request handling remains outside StorageService.
def upload_audio_route(file, db):
    stored_path = StorageService().save_uploaded_wav(file)
    return create_uploaded_audio_record(db, file.filename, stored_path)
\end{lstlisting}

In this example, the storage class does not create separation jobs and does not run NeoSSNet. Its responsibility is limited to saving files and returning a path that other components can use.

\subsection{Open/Closed Principle}

The Open/Closed Principle suggests that the system should be open for extension but closed for unnecessary modification. The separation workflow should be able to support a new implemented algorithm by adding a strategy class and registering it in the factory, without rewriting upload, result, history, or download routes.

\begin{lstlisting}[caption={Open/Closed Principle Example}]
# Open for extension:
# add a new implemented strategy and register it in the factory.
class NeoSSNetStrategy:
    def separate(self, input_wav_path, model_path, config_path,
                 heart_output_path, lung_output_path):
        return run_neossnet_inference(...)


class SeparationAlgorithmFactory:
    _registry = {
        "neossnet": NeoSSNetStrategy,
    }

    @classmethod
    def create(cls, model):
        strategy_class = cls._registry[model.architecture.lower()]
        return strategy_class()
\end{lstlisting}

The workflow is closed to unnecessary modification because \texttt{SeparationService} can keep the same high-level steps. A future algorithm should mainly require a new concrete strategy and a factory registration.

\subsection{Dependency Inversion Principle}

Dependency Inversion is relevant because high-level workflow classes should not be tightly bound to one concrete model implementation. \texttt{SeparationService} should depend on the \texttt{SeparationAlgorithm} abstraction, while \texttt{NeoSSNetStrategy} provides the current concrete implementation.

\begin{lstlisting}[caption={Dependency Inversion Principle Example}]
# Dependency Inversion Principle:
# high-level workflow depends on the abstraction.
class SeparationService:
    def run_algorithm(self, algorithm: SeparationAlgorithm, paths):
        return algorithm.separate(
            input_wav_path=paths.input_wav,
            model_path=paths.model_checkpoint,
            model_config_path=paths.model_config,
            heart_output_path=paths.heart_output,
            lung_output_path=paths.lung_output,
        )
\end{lstlisting}

The service does not need to instantiate or know the internal details of \texttt{NeoSSNetStrategy}. The factory supplies an object that satisfies the \texttt{SeparationAlgorithm} interface, and the service runs the workflow through that abstraction.

\subsection{Interface-Based Design for Algorithm Replacement}

Interface-based design supports future algorithm replacement by defining the expected input and output of a separation algorithm. The current algorithm should produce separated heart and lung audio outputs, but the service layer should not need to know whether the implementation uses NeoSSNet, an ONNX-exported model, or a future baseline method. For the current project, NeoSSNet remains the required real separation model.

\subsection{Persistence Separation Between SQLite and Filesystem}

The design separates structured metadata from large binary audio files. SQLite stores records such as filenames, file paths, upload timestamps, job status, model information, output paths, metrics, and logs. The filesystem stores uploaded WAV files, generated heart and lung outputs, model checkpoints, and dataset files. This approach is simpler to run locally and avoids making the database unnecessarily large.

\section{Software Design Plan}

The proposed software design is organised into layers. This layered plan supports a clear implementation path and makes the system easier to explain in the report and viva presentation.

\begin{table}[H]
  \centering
  \small
  \caption{Software Design Plan by Layer}
  \label{tab:software-design-plan}
  \begin{tabularx}{\textwidth}{p{0.24\textwidth} p{0.30\textwidth} Y}
    \toprule
    \textbf{Layer} & \textbf{Main Responsibility} & \textbf{Project Role} \\
    \midrule
    Frontend layer & HTML, CSS, JavaScript, templates, and browser audio controls & Provides upload form, status messages, original and separated audio playback, download links, and history display. \\
    Backend API layer & FastAPI application and routers & Receives upload, separation, result, download, history, log, and health requests. \\
    Service layer & Workflow coordination and business logic & Connects validation, database updates, NeoSSNet inference, output storage, logging, and result retrieval. \\
    ML inference layer & Audio preprocessing and model inference & Loads or prepares audio, applies the NeoSSNet model, and returns separated heart and lung waveforms. \\
    Database layer & SQLite schema and repository-style access & Stores uploaded audio metadata, model records, job states, result paths, metrics, and system logs. \\
    Storage layer & Local filesystem folders & Stores uploaded mixed WAV files, temporary files, separated heart and lung outputs, logs where needed, and ML model weights. \\
    \bottomrule
  \end{tabularx}
\end{table}

The design plan follows the project priority of real functionality first, then simplicity and local execution. It avoids unnecessary cloud, microservice, or authentication complexity unless those features are explicitly added later.
